% s-366-p366-m002; original PDF 366, 367
\Aphorism{s-366-p366-m002}{409}
Die erkenntnisstheoretische Hauptschwierigkeit, wie ein transcendent willkürliches Geschehen empirisch als gesetzlich geregelter Weltverlauf „erscheinen“ könne, verschwindet bei folgender Betrachtung. Einzig real ist die Gegenwart, genauer das Zeitdifferential, das zum Zustandekommen eines Elementarbewusstseins mindestens erforderlich ist (die Gegenwart, als Zeitpunkt gedacht, ist noch nicht Bewusstseinsphänomen); alles Weitere, d. h. die ganze Vergangenheit und Zukunft, ist ungeheure Extrapolation, inhaltliche Weiterverlängerung des Zeitdifferentials zur Zeitlinie, aber der Existenz nach unbeweisbar. Man kann sich darüber in populärer Weise klar werden, wenn man die unmittelbare Gegenwart irgendwie, etwa durch Aussprechen des Wortes ‚jetzt‘ zu markiren sucht; in dem Moment, da wir den Endconsonanten t aussprechen, ist der Augenblick der phonetischen Erzeugung des j bereits Vergangenheit, d. h. Glaubenssache: eine perspectivische Annahme, gegründet auf inhaltlichen regressus von der Gegenwart aus. Wenn wir also behaupten, in einer gesetzlich zusammenhängenden, qualitativ einheitlichen, unser causales Bedürfniss vollkommen befriedigenden Welt zu leben (dies ist die postulirte, wenn auch noch nicht practisch durchgeführte Anschauung der Naturwissenschaft), so bedeutet dies zunächst nur die Illusion eines Zeitdifferentials, kürzer als die Zeitdauer des gesprochenen Wortes jetzt; eine Illusion aber, zu deren Art es gehört, sich als Nachfolgerin und Vorgängerin einer gleichartigen Illusion zu empfinden. Es giebt — das wissen wir — Zeitelemente, die derartig empfinden; dass es nur solche gebe, liegt im Wesen der Illusion, aber nicht der Sache. Die Illusion scheint seit Ewigkeit bestanden zu haben und in Ewigkeit zu bestehen; aber eben dieser Schein ist nur ein Theil der Illusion, die darum trotzdem jeden Augenblick durchbrochen werden kann.
\par

